% Chapter 2: Requirement Specification
% Main chapter title is set in main.tex via \chapter{Requirement Specification}

\section{Specific Requirements}

This section describes the functional and non-functional requirements of the Hybrid Post-Quantum VPN system. These requirements define the operational capabilities, security properties, performance expectations, and deployment constraints of the proposed hybrid VPN architecture.

\subsection{Functional Requirements}

Functional requirements describe the specific operations and services that the system must perform.

\subsubsection{Hybrid Handshake Module}

The hybrid handshake module is responsible for establishing secure session keys using both classical and post-quantum cryptographic primitives.

\begin{itemize}

    \item The system shall perform X25519 elliptic-curve Diffie-Hellman key exchange during every tunnel handshake.

    \item The system shall perform ML-KEM-768 post-quantum key encapsulation in parallel with the classical key exchange.

    \item The system shall combine the X25519 and ML-KEM-768 shared secrets using HKDF-SHA256 to derive final session keys.

    \item The system shall bind derived session keys to the handshake transcript to prevent downgrade attacks and transcript tampering.

    \item The system shall support secure algorithm negotiation between communicating peers.

    \item The system shall authenticate peers using ECDSA-P256 signatures over the handshake transcript.

    \item The system shall support future integration of ML-DSA hybrid post-quantum signatures.

\end{itemize}

\subsubsection{Tunnel and Encryption Module}

The tunnel module is responsible for encrypted communication and secure packet forwarding.

\begin{itemize}

    \item The system shall encrypt all tunnel traffic using AES-256-GCM authenticated encryption.

    \item The system shall include a unique nonce and authentication tag for every encrypted packet.

    \item The system shall implement replay protection using packet sequence numbers and nonce verification.

    \item The system shall drop and log packets that fail integrity verification.

    \item The system shall securely forward IPv4 packets between Linux TUN interfaces and encrypted UDP sockets.

    \item The system shall support bidirectional encrypted communication between client and gateway nodes.

    \item The system shall trigger rekeying upon nonce exhaustion or prolonged tunnel usage.

\end{itemize}

\subsubsection{Network Interface Module}

The network interface module handles virtual networking and routing operations.

\begin{itemize}

    \item The system shall create Linux TUN interfaces dynamically upon successful tunnel establishment.

    \item The system shall configure IP addresses and routing rules using \texttt{pyroute2}.

    \item The system shall support encrypted UDP transport over configurable ports.

    \item The system shall cleanly tear down TUN interfaces and routes after tunnel termination.

    \item The system shall support communication between virtual machines connected through bridged or NAT networking.

\end{itemize}

\subsubsection{Desktop Dashboard Module}

The desktop dashboard provides monitoring and control capabilities for the VPN tunnel.

\begin{itemize}

    \item The system shall provide an Electron + Next.js desktop dashboard for tunnel management.

    \item The dashboard shall allow users to initiate and terminate VPN connections.

    \item The dashboard shall display real-time tunnel status including connection state, active algorithms, throughput, and session details.

    \item The dashboard shall display handshake logs and tunnel events.

    \item The system shall expose REST APIs for frontend-backend communication using FastAPI.

    \item The frontend shall support browser-only development mode for testing purposes.

\end{itemize}

\subsubsection{Observability and Testing Module}

The observability module provides monitoring, logging, and experimental evaluation capabilities.

\begin{itemize}

    \item The system shall maintain timestamped logs of handshake and tunnel events.

    \item The system shall expose health-check and status APIs for monitoring.

    \item The system shall provide metrics for handshake latency, throughput, packet counts, and replay detection.

    \item The system shall support experimental comparison between classical-only and hybrid handshakes.

    \item The system shall provide automated unit testing using \texttt{pytest}.

\end{itemize}

\subsection{Non-Functional Requirements}

Non-functional requirements define the quality attributes and operational constraints of the system.

\begin{itemize}

    \item \textbf{Performance:} The hybrid handshake shall introduce minimal latency overhead compared to classical-only VPN handshakes. The encrypted tunnel shall sustain practical throughput for secure remote communication and file transfer.

    \item \textbf{Security:} The system shall provide confidentiality, integrity, replay protection, and forward secrecy. Session keys shall depend on both classical and post-quantum cryptographic mechanisms.

    \item \textbf{Reliability:} The system shall detect tunnel disconnections and malformed packets gracefully without crashing or entering undefined states.

    \item \textbf{Usability:} The Electron + Next.js dashboard shall provide an intuitive user interface for tunnel monitoring and management.

    \item \textbf{Maintainability:} Cryptographic primitives shall be modular and abstracted to support future replacement or extension with alternative post-quantum algorithms.

    \item \textbf{Scalability:} The architecture shall support future extension toward enterprise-scale VPN deployments and additional cryptographic suites.

    \item \textbf{Deployability:} The system shall support deployment using Docker containers and Linux virtual machines.

\end{itemize}

\section{Hardware Requirements}

The Hybrid Post-Quantum VPN system requires Linux-compatible hardware capable of supporting encrypted tunneling, virtual networking, and post-quantum cryptographic operations.

\begin{table}[H]
\centering
\begin{tabular}{|l|p{8cm}|}
\hline
\textbf{Component} & \textbf{Specification} \\ \hline

Processor & Quad-core x86\_64 CPU with AES-NI support (8-core with AVX2 recommended) \\ \hline

RAM & 8 GB DDR4 minimum (16 GB recommended) \\ \hline

Storage & 40 GB SSD minimum (100 GB NVMe SSD recommended) \\ \hline

Network Interface & 1 Gbps Ethernet connection (wired preferred for latency testing) \\ \hline

Operating System & Ubuntu 22.04/24.04 LTS with Linux kernel version 5.15 or higher \\ \hline

Kernel Support & Linux TUN/TAP driver enabled using \texttt{modprobe tun} \\ \hline

\end{tabular}
\caption{Hardware Requirements}
\end{table}

\section{Software Requirements}

The system uses modern networking, cryptographic, backend, and frontend technologies to implement a research-grade hybrid VPN prototype.

\subsection{Development Environment}

\begin{table}[H]
\centering
\begin{tabular}{|l|l|l|}
\hline
\textbf{Category} & \textbf{Technology} & \textbf{Version} \\ \hline

Operating System & Ubuntu Linux & 22.04/24.04 LTS \\ \hline

Programming Language & Python & 3.11 or higher \\ \hline

Frontend Runtime & Node.js & 20 LTS or higher \\ \hline

Package Manager & pnpm / uv & Latest \\ \hline

Version Control & Git & 2.40 or higher \\ \hline

IDE & Visual Studio Code & Latest \\ \hline

\end{tabular}
\caption{Development Environment}
\end{table}

\subsection{Application Stack}

\begin{table}[H]
\centering
\begin{tabular}{|l|l|l|p{4cm}|}
\hline
\textbf{Layer} & \textbf{Technology} & \textbf{Version} & \textbf{Purpose} \\ \hline

Frontend Framework & Next.js & 14+ & Desktop dashboard and UI \\ \hline

Desktop Runtime & Electron & Latest & Cross-platform desktop application \\ \hline

Frontend Library & React & Latest & User interface rendering \\ \hline

Styling & Tailwind CSS & 3.4+ & Responsive UI styling \\ \hline

Frontend Language & TypeScript & 5+ & Type-safe frontend development \\ \hline

Backend Framework & FastAPI & 0.110+ & REST API services \\ \hline

Backend Runtime & Python & 3.11+ & Backend implementation \\ \hline

Async Runtime & AsyncIO & Built-in & Concurrent networking operations \\ \hline

Cryptographic Library & OpenSSL & 3.0+ & X25519, AES-GCM, HKDF, ECDSA \\ \hline

Post-Quantum Library & liboqs & 0.10+ & ML-KEM-768 and ML-DSA \\ \hline

Python Bindings & pyoqs & 0.10+ & Python interface for liboqs \\ \hline

Networking & pyroute2 & Latest & Linux TUN/TAP and route management \\ \hline

Transport Protocol & UDP Sockets & Native & Encrypted tunnel transport \\ \hline

Testing Framework & pytest & 8+ & Automated backend testing \\ \hline

Containerization & Docker & Latest & Deployment and testing \\ \hline

\end{tabular}
\caption{Application Stack}
\end{table}

\subsection{Deployment}

\begin{table}[H]
\centering
\begin{tabular}{|l|l|p{5cm}|}
\hline
\textbf{Component} & \textbf{Technology} & \textbf{Purpose} \\ \hline

Backend Deployment & Docker & Containerized backend deployment \\ \hline

Frontend Deployment & Electron Desktop App & Desktop-based VPN controller \\ \hline

Networking Environment & Linux Virtual Machines & Tunnel and routing evaluation \\ \hline

Hypervisor & VMware Workstation / VirtualBox & VM-based testing \\ \hline

Monitoring Tools & Wireshark / tcpdump & Packet inspection and debugging \\ \hline

API Hosting & FastAPI + Uvicorn & Backend API services \\ \hline

Transport Layer & UDP over IPv4 & Encrypted tunnel communication \\ \hline

\end{tabular}
\caption{Deployment Stack}
\end{table}